
\documentclass[addpoints]{exam}
\usepackage{amsmath}
\usepackage{amssymb}
\usepackage{color}

% \printanswers

\newcommand{\event}{Limits and Rates of Change}
\newcommand{\tournament}{}
\def\type{0}

\setlength\fillinlinelength{\textwidth}
\CorrectChoiceEmphasis{\color{red}\bfseries}
\SolutionEmphasis{\color{red}\bfseries}

\setlength\answerclearance{2pt}
\pagestyle{headandfoot}
\runningheadrule
\runningheader{\event}{\tournament}{\ifnum\type=1 Class Set Test \else Full Name:\hspace{1cm} \fi}
\runningfootrule
\runningfooter{}{Page \thepage\ of \numpages}{}

\begin{document}
\begin{center}
    {\huge Grade 12 Calculus \& Vectors \\ \event
    \ifprintanswers \space Key
    \else \space \ifnum\type<2 Test \fi \ifnum\type=1 \textbf{(CLASS SET)} \fi
          \ifnum\type=2 Answer Sheet \fi
    \fi \\ \vspace{3mm}\tournament\vspace{1mm}}
\end{center}

\vspace{.25\textheight}

\begin{center}
    First Name: \ifprintanswers
    \underline{\hspace{3.5cm}}\textbf{KEY}\underline{\hspace{5.5cm}}\\\vspace{5mm}
    \else \underline{\hspace{10cm}}\\ \vspace{5mm} \fi
    Last Name: \underline{\hspace{10cm}}\\ \vspace{5mm}
\end{center}

\noindent\textbf{\underline{Directions:}}
\begin{itemize}
    \ifnum\type=1 \item \textbf{This test is a class set.} Do not write on this paper. \fi
    \item Show all work for full marks on short answer questions.
    \item The test is designed to be completed in 75 minutes.
\end{itemize}
\begin{center}
\textit{For grading use only}\\
\multirowgradetable{1}[pages]
\end{center}

\newpage
\begin{questions}

\section*{Part A --- Multiple Choice (15 marks)}
\vspace{5mm}

\question[1] The average rate of change of $f(x) = x^2$ from $x = 1$ to $x = 3$ is:
\begin{choices}
    \choice 2
    \correctchoice 4
    \choice 6
    \choice 8
\end{choices}

\question[1] $\displaystyle\lim_{x \to 2} \frac{x^2 - 4}{x - 2}$ equals:
\begin{choices}
    \choice 0
    \choice 2
    \correctchoice 4
    \choice Does not exist
\end{choices}

\question[1] $\displaystyle\lim_{x \to 0} \frac{\sin x}{x}$ equals:
\begin{choices}
    \choice 0
    \correctchoice 1
    \choice $\infty$
    \choice Does not exist
\end{choices}

\question[1] $\displaystyle\lim_{x \to \infty} \frac{5x^2 - 3x}{2x^2 + 1}$ equals:
\begin{choices}
    \choice 0
    \choice $-3$
    \correctchoice $\dfrac{5}{2}$
    \choice $\infty$
\end{choices}

\question[1] $\displaystyle\lim_{x \to 3} \frac{x^2 - 9}{x - 3}$ equals:
\begin{choices}
    \choice 0
    \choice 3
    \correctchoice 6
    \choice Does not exist
\end{choices}

\question[1] Which expression is the correct definition of $f'(a)$?
\begin{choices}
    \choice $\displaystyle\lim_{h \to 0} \frac{f(a+h) + f(a)}{h}$
    \correctchoice $\displaystyle\lim_{h \to 0} \frac{f(a+h) - f(a)}{h}$
    \choice $\displaystyle\lim_{h \to 0} \frac{f(a) - f(h)}{h}$
    \choice $\displaystyle\lim_{h \to 0} \left[f(a+h) - f(a)\right] \cdot h$
\end{choices}

\question[1] A function $f(x)$ is continuous at $x = c$ if and only if:
\begin{choices}
    \choice $f(c)$ is defined
    \choice $\displaystyle\lim_{x \to c} f(x)$ exists
    \correctchoice $\displaystyle\lim_{x \to c} f(x) = f(c)$
    \choice $f(c) = 0$
\end{choices}

\question[1] $\displaystyle\lim_{x \to 0} \frac{\sqrt{x+4} - 2}{x}$ equals:
\begin{choices}
    \choice 0
    \correctchoice $\dfrac{1}{4}$
    \choice 1
    \choice Does not exist
\end{choices}

\question[1] Which function is \textbf{not} continuous at $x = 2$?
\begin{choices}
    \choice $f(x) = x^2 - 4$
    \correctchoice $f(x) = \dfrac{x^2 - 4}{x - 2}$
    \choice $f(x) = |x - 2|$
    \choice $f(x) = 3$
\end{choices}

\question[1] $\displaystyle\lim_{x \to 4} \frac{x - 4}{\sqrt{x} - 2}$ equals:
\begin{choices}
    \choice 0
    \choice 2
    \correctchoice 4
    \choice 8
\end{choices}

\question[1] The slope of the tangent to $f(x) = x^2$ at the point $(3,\, 9)$ is:
\begin{choices}
    \choice 3
    \correctchoice 6
    \choice 9
    \choice 18
\end{choices}

\question[1] For $f(x) = |x|$, which statement is true?
\begin{choices}
    \choice $f$ is differentiable everywhere.
    \choice $f$ is not continuous at $x = 0$.
    \correctchoice $f$ is continuous but not differentiable at $x = 0$.
    \choice $f$ is neither continuous nor differentiable at $x = 0$.
\end{choices}

\question[1] $\displaystyle\lim_{x \to -\infty} \frac{2x - 1}{x + 3}$ equals:
\begin{choices}
    \choice $-\dfrac{1}{3}$
    \choice $-2$
    \correctchoice $2$
    \choice Does not exist
\end{choices}

\question[1] The average velocity of a particle with position $s(t) = t^2 + 2t$
from $t = 1$ to $t = 3$ is:
\begin{choices}
    \choice 4 m/s
    \correctchoice 6 m/s
    \choice 8 m/s
    \choice 10 m/s
\end{choices}

\question[1] The instantaneous rate of change of $f(x) = x^3$ at $x = 2$ is:
\begin{choices}
    \choice 4
    \choice 8
    \correctchoice 12
    \choice 24
\end{choices}


\section*{Part B --- Short Answer (33 marks)}

\question The position of a particle moving along a line is given by
$s(t) = 2t^2 - 3t + 1$ (metres, $t$ in seconds).
\begin{parts}
    \part[2] Calculate the average velocity from $t = 1$ to $t = 4$.
    \part[3] Write a simplified expression for the average velocity from $t = 1$ to $t = 1 + h$.
    \part[2] Use your answer to find the instantaneous velocity at $t = 1$.
    \part[1] Interpret the sign of the instantaneous velocity at $t = 1$ in context.
\end{parts}
\begin{solution}[9cm]
\textbf{(a)} $s(1) = 2 - 3 + 1 = 0$;\quad $s(4) = 32 - 12 + 1 = 21$.
\[
\text{Avg vel} = \frac{s(4) - s(1)}{4 - 1} = \frac{21 - 0}{3} = \mathbf{7 \text{ m/s}}
\]

\textbf{(b)}
$s(1+h) = 2(1+h)^2 - 3(1+h) + 1 = 2 + 4h + 2h^2 - 3 - 3h + 1 = h + 2h^2$
\[
\text{Avg vel} = \frac{s(1+h) - s(1)}{h} = \frac{h + 2h^2}{h} = \mathbf{1 + 2h}
\]

\textbf{(c)}
\[
\lim_{h \to 0}(1 + 2h) = \mathbf{1 \text{ m/s}}
\]

\textbf{(d)} The particle is moving in the positive direction at 1 m/s at $t = 1$ s.
\end{solution}

\question Evaluate each limit. Show all algebraic steps; direct substitution alone earns no marks
where it produces an indeterminate form.
\begin{parts}
    \part[2] $\displaystyle\lim_{x \to 3} \frac{x^2 - 9}{x - 3}$
    \part[2] $\displaystyle\lim_{x \to 0} \frac{\sqrt{x + 4} - 2}{x}$
    \part[2] $\displaystyle\lim_{x \to \infty} \frac{3x^2 - 2x + 1}{x^2 + 5}$
    \part[2] $\displaystyle\lim_{x \to 2^-} f(x)$, where
    $f(x) = \begin{cases} x + 1 & x < 2 \\ x^2 - 1 & x \geq 2 \end{cases}$
\end{parts}
\begin{solution}[9cm]
\textbf{(a)} Factor: $\dfrac{(x+3)(x-3)}{x-3} = x + 3 \to \mathbf{6}$

\textbf{(b)} Rationalize by multiplying by $\dfrac{\sqrt{x+4}+2}{\sqrt{x+4}+2}$:
\[
\frac{(x+4) - 4}{x(\sqrt{x+4}+2)} = \frac{x}{x(\sqrt{x+4}+2)} = \frac{1}{\sqrt{x+4}+2}
\to \frac{1}{2+2} = \mathbf{\frac{1}{4}}
\]

\textbf{(c)} Divide numerator and denominator by $x^2$:
$\dfrac{3 - 2/x + 1/x^2}{1 + 5/x^2} \to \mathbf{3}$

\textbf{(d)} As $x \to 2^-$, use $f(x) = x + 1$: $\quad \lim = 2 + 1 = \mathbf{3}$
\end{solution}

\question Using the \textbf{definition of the derivative} (first principles), find $f'(x)$ for
$f(x) = 3x^2 - 2x$. Show all steps fully.
\begin{parts}
    \part[5] Derive $f'(x)$ from first principles.
    \part[3] Use your result to find the equation of the tangent line to $f$ at $x = 1$.
\end{parts}
\begin{solution}[9cm]
\textbf{(a)}
\begin{align*}
f'(x) &= \lim_{h \to 0} \frac{f(x+h) - f(x)}{h} \\
&= \lim_{h \to 0} \frac{[3(x+h)^2 - 2(x+h)] - [3x^2 - 2x]}{h} \\
&= \lim_{h \to 0} \frac{3x^2 + 6xh + 3h^2 - 2x - 2h - 3x^2 + 2x}{h} \\
&= \lim_{h \to 0} \frac{6xh + 3h^2 - 2h}{h} \\
&= \lim_{h \to 0} (6x + 3h - 2) = \boxed{6x - 2}
\end{align*}

\textbf{(b)} $f(1) = 3 - 2 = 1$;\quad $f'(1) = 6 - 2 = 4$.
\[
y - 1 = 4(x - 1) \implies \mathbf{y = 4x - 3}
\]
\end{solution}

\question A function is defined as:
\[
f(x) = \begin{cases} x^2 - 1 & x < 2 \\ 3 & x = 2 \\ 2x - 1 & x > 2 \end{cases}
\]
\begin{parts}
    \part[1] Find $\displaystyle\lim_{x \to 2^-} f(x)$.
    \part[1] Find $\displaystyle\lim_{x \to 2^+} f(x)$.
    \part[2] Does $\displaystyle\lim_{x \to 2} f(x)$ exist? Justify your answer.
    \part[2] Is $f$ continuous at $x = 2$? Justify using the formal definition of continuity.
    \part[2] Sketch a graph of $f$ near $x = 2$, indicating any discontinuities or holes.
\end{parts}
\begin{solution}[8cm]
\textbf{(a)} $\displaystyle\lim_{x \to 2^-}(x^2 - 1) = 4 - 1 = \mathbf{3}$

\textbf{(b)} $\displaystyle\lim_{x \to 2^+}(2x - 1) = 4 - 1 = \mathbf{3}$

\textbf{(c)} Yes. Since both one-sided limits equal 3, $\displaystyle\lim_{x \to 2} f(x) = 3$.

\textbf{(d)} Yes, $f$ is continuous at $x = 2$. All three conditions hold:
(i) $f(2) = 3$ is defined;\quad (ii) $\displaystyle\lim_{x \to 2}f(x) = 3$ exists;\quad
(iii) $\displaystyle\lim_{x \to 2}f(x) = f(2) = 3$. \checkmark

\textbf{(e)} The graph is a parabola for $x < 2$, a filled dot at $(2,3)$, and a line for $x > 2$.
All three pieces connect at the single point $(2, 3)$ --- no break or hole.
\end{solution}

\end{questions}
\end{document}
